\section{Graduating to L1}
\label{sec:l1-graduation}

Through 2.x, \Zoo{} ran as a \Lux{} L2 application: an \textsc{evm}-compatible
chain (Chain ID 200200) settling state commitments to \Lux{}'s C-Chain. With
3.0 \Zoo{} graduates to a sovereign L1 with its own validator set, its own
\Quasar{}Cert finality, and its own economic security parameters, while
remaining inside the \Lux{} chain topology.

\subsection{What changes}

\begin{longtable}{p{0.30\textwidth}p{0.30\textwidth}p{0.34\textwidth}}
\toprule
\textbf{Concern} & \textbf{Zoo 2.x (L2)} & \textbf{Zoo 3.0 (L1)} \\
\midrule
Validator set & inherited from Lux & sovereign Zoo validators \\
Finality & via Lux C-Chain & \Quasar{}Cert local finality \\
Economic security & rented from Lux & \$ZOO + \$AI staking \\
Slashing & N/A & enforced by Q-Chain ceremony \\
Block production & sequencer & full BFT replication \\
Bridge to Lux & state commitments & B-Chain (\textsc{lp-134}) \\
Settlement guarantees & rollup proof & sovereign certificate \\
Native gas asset & ZOO (bridged) & ZOO (native) \\
\bottomrule
\end{longtable}

\subsection{Why graduate}

\paragraph{Determinism for AI workloads.} Per-\LLM{} chain execution
\cite{zoo-per-llm-chains} requires deterministic scheduling of training-data
attestations, attribution graph updates, and inference receipt rollups.
Running these as L2 transactions inside a generic L1 sequencer was a
solvable problem in 2.x; running them as native operations on a sovereign
chain whose mempool, consensus, and finality are tuned for AI is materially
better.

\paragraph{Native confidential compute.} \Zoo{} 3.0 exposes
F-Chain-backed \FHE{} (\textsc{lp-013}) and Z-Chain-backed \ZKP{}
(\textsc{lp-063}) as first-class precompiles. As an L2, every confidential
operation paid the cross-domain message tax. As an L1 with chain-local
precompiles, the same operations are a single transaction.

\paragraph{Sovereign tokenomics.} Zoo's PoAI ledger and AI mining rewards
need a chain where the validator subsidy schedule, fee burn policy, and
\$ZOO/\$AI dual-asset gas semantics are governance variables, not L2
constants inherited from upstream. L1 graduation puts these levers where
the Zoo \DAO{} can move them.

\subsection{Constructive embedding in Lux topology}

L1 graduation is not departure. The \Quasar{}-Native App Stack
(\textsc{lp-133}) defines exactly how a sovereign appchain settles into the
broader \Lux{} topology (\textsc{lp-134}):

\begin{itemize}[leftmargin=1.2em]
  \item \textbf{Q-Chain} runs the \Corona{} threshold ceremonies that
        validate \Zoo{}'s validator key material.
  \item \textbf{A-Chain} (AIVM) anchors TEE quotes for AI mining receipts;
        \Zoo{} validators submit attestations there.
  \item \textbf{B-Chain} (BVM) carries \Zoo{}$\leftrightarrow$\Lux{} bridge
        messages and \Zoo{}$\leftrightarrow$external-chain transfers via
        \Zoo{} Bridge (\S\ref{sec:zoo-bridge}).
  \item \textbf{F-Chain} (FVM) and \textbf{Z-Chain} (ZVM) provide
        confidential compute substrate that \Zoo{} dispatches to via
        precompiles.
\end{itemize}

This is the appchain pattern: sovereign block production, shared cross-chain
plumbing.

\subsection{Validator economics}

Validators stake \$ZOO and optionally \$AI. Block rewards are paid in the
gas asset of the block. Quasar 3.0 triple-consensus
(\textsc{lp-020}: BLS + Corona + ML-DSA) means every signed block carries
both a classical fast path proof and a post-quantum proof; validators
running compromised classical keys still cannot forge blocks because the
\Corona{} layer must independently authorize. Slashing is enforced via
Q-Chain ceremonies for double-sign and via A-Chain attestation gaps for
liveness failures.
